\documentclass[12pt]{article}

\setlength{\parindent}{0pt}
\setcounter{secnumdepth}{3}
\usepackage[margin=1in]{geometry}

% packages
\usepackage{amsmath, amsfonts, amsthm, amssymb}
\usepackage{graphicx, float}
\usepackage{mathtools}
\usepackage{titlesec}
\usepackage{interval}
\usepackage{hyperref}
\usepackage{siunitx}
\usepackage{titling}
\usepackage{vwcol}
\usepackage{setspace}
\usepackage{empheq}
\usepackage{cancel}
\usepackage{esdiff}
\usepackage{multicol}
\usepackage{mdframed}
\usepackage{esdiff}
\usepackage{tikzsymbols}
\usepackage{multicol}
\usepackage{tikz}
\usepackage{varwidth}
\usepackage{pgfplots}
\usepackage{fancyhdr}
\usepackage{enumitem}
\usepackage{tcolorbox}

\intervalconfig {
	soft open fences
}

% header
\pagestyle{fancy}
\fancyhf{}
\lhead{Ethan Fan}
\rhead{AP Calculus BC}
\cfoot{\thepage}

% section format
\titleformat{\section}
  {\large \bfseries}
  {}
  {0em}
  {}
\titlespacing*{\section}{0pt}{0.5em}{0.3em}

\titleformat{\subsection}[runin]
  {\bfseries}
  {}
  {0em}
  {}
\titlespacing*{\subsection}{0pt}{0pt}{0em}

\titleformat{\subsubsection}[runin]
  {\itshape}
  {(\roman{subsubsection})}
  {0em}
  {}

% commands
\newcommand{\alignedintertext}[1]{%
  \noalign{%
    \vskip\belowdisplayshortskip
    \vtop{\hsize=\linewidth#1\par
    \expandafter}%
    \expandafter\prevdepth\the\prevdepth
  }%
}

\newcommand*{\paren}[1]{\ensuremath\left(#1\right)}
\newcommand*{\limit}[2][x]{\ensuremath{\displaystyle\lim_{#1\to#2}}}
\newcommand*{\Limit}[3][x]{\ensuremath{\displaystyle\lim_{#1\to#2}\left[#3\right]}}
\newcommand*{\deriv}[1][x]{\ensuremath{\dfrac{\mathrm{d}}{\mathrm{d}#1}}}
\newcommand*{\Deriv}[2][x]{\ensuremath{\dfrac{\mathrm{d}}{\mathrm{d}#1}\left[#2\right]}}
\newcommand{\ddx}{\dfrac{d}{dx}}
\newcommand{\dydx}{\dfrac{dy}{dx}}
\newcommand{\dydt}{\dfrac{dy}{dt}}
\newcommand{\dxdt}{\dfrac{dx}{dt}}
\newcommand*{\iinteg}[2][x]{\ensuremath{\displaystyle\int #2\;\mathrm{d}#1}}
\newcommand*{\dinteg}[4][x]{\ensuremath{\displaystyle\int_{#2}^{#3}#4\;\mathrm{d}#1}}
\newcommand*{\abs}[1]{\ensuremath{\left|#1\right|}}
\newcommand*{\eps}{\varepsilon}
\newcommand*{\real}{\ensuremath\mathbb{R}}
\newcommand*{\integer}{\ensuremath\mathbb{Z}}
\newcommand*{\posint}{\ensuremath\mathbb{Z^+}}
\newcommand*{\cbrt}[1]{\ensuremath{\sqrt[3]{#1}}}
\newcommand*{\lheqzero}{\ensuremath{\underset{\text{L'H}}{\overset{\left[\frac00\right]}{=}}}}
\newcommand*{\lheqinfty}{\ensuremath{\underset{\text{L'H}}{\overset{\left[\frac{\infty}{\infty}\right]}{=}}}}
\newcommand{\tor}{\text{ or }}
\newcommand{\floor}[1]{\lfloor #1 \rfloor}

\newcommand{\vem}{\vspace{0.5em}}
\newcommand{\example}[1]{\section{Example #1}}
\newcommand{\problem}[1]{\section{Problem #1}}
\newcommand{\subproblem}[1]{\subsection{(#1)} \,}

\DeclareMathOperator{\dne}{DNE}
\DeclareMathOperator{\sgn}{sgn}

\newcommand{\definition}[1]{\begin{tcolorbox}[colback=red!5!white,colframe=red!75!black,parbox=false] #1 \end{tcolorbox}}

\newcommand{\theorem}[2]{\begin{tcolorbox}[title={#1},colback=blue!5!white,colframe=blue!75!black,parbox=false] #2 \end{tcolorbox}}

\allowdisplaybreaks
% \postdisplaypenalty=100000

% document
\begin{document}

\vspace*{0cm}

% opening
\begin{center}
    {\Large \bfseries Problem Set 5} \\[0.5cm]
    {\large Antiderivatives} \\[0.35cm]
    {\normalsize September 15, 2026}
\end{center}

% problems

\problem{3}

\subproblem{j} 
\begin{gather*}
    \int \cos^3 x \, dx\\
    =\int \frac{\cos 3x}{4}+\frac{3\cos x}{4} \, dx\\
    =\boxed{\frac{\sin 3x}{12}+\frac{3\sin x}{4}+C}
\end{gather*}

\subproblem{k}
\begin{gather*}
    \int \frac{(x^3+8)(x-1)}{x^2-2x+4} \, dx\\
    =\int (x+2)(x-1) \, dx\\
    =\int x^2+x-2 \, dx\\
    =\boxed{\frac{x^3}{3}+\frac{x^2}{2}-2x+C}
\end{gather*}

\subproblem{l}
\begin{gather*}
    \int \frac{1}{\sqrt{x+1}-\sqrt{x-1}} \, dx\\
    =\int \frac{\sqrt{x+1}+\sqrt{x-1}}{2} \, dx\\
    =\frac{1}{2}\cdot \frac{2}{3} \cdot ((x+1)^\frac{3}{2}+(x-1)^\frac{3}{2})+C\\
    =\boxed{\frac{1}{3}(x+1)^\frac{3}{2}+\frac{1}{3}(x-1)^\frac{3}{2}+C}
\end{gather*}

\subproblem{m}
\begin{gather*}
    \int \sin^4 x \, dx\\
    =\int \paren{\frac{1-\cos 2x}{2}}^2 \, dx\\
    =\frac{1}{4}\int \cos^2(2x)-2\cos 2x+1\, dx\\
    =\frac{1}{4}\int \frac{\cos 4x+1}{2}-2\cos 2x+1 \, dx\\
    =\int \frac{3}{8}  - \frac{\cos 2x}{2} + \frac{\cos 4x}{8} \, dx\\
    =\boxed{\frac{3}{8}x- \frac{\sin 2x}{4}+ \frac{\sin 4x}{32}+C}
\end{gather*}

\subproblem{n}
\begin{gather*}
    \int \sin x \cos x \cos 2x \cos 4x \cos 8x \, dx\\
    =\int \frac{\sin 16x}{16} \, dx\\
    =\boxed{-\frac{\cos 16x}{256}+C}
\end{gather*}

\problem{4}

\subproblem{a}
\begin{gather*}
    f'(x)=6\sqrt{x}+5x^\frac{3}{2}\\
    f(x)=4x^\frac{3}{2}+2x^\frac{5}{2}+C\\
    f(1)=10 \implies c=4\\
    \boxed{f(x)=4x^\frac{3}{2}+2x^\frac{5}{2}+4}
\end{gather*}

\subproblem{b}
\begin{gather*}
    f''(\theta)=\sin \theta + \cos \theta \\
    f'(\theta )= \sin \theta -\cos \theta +C\\
    f'(0)=1 \implies C=2\\
    f(\theta )=-\cos \theta -\sin \theta +2\theta +C\\
    f(0)=2 \implies C=3\\
    \boxed{f(\theta )=-\cos \theta -\sin \theta +2\theta +3}
\end{gather*}
Two conditions are needed, as we are given the second derivative. Similarly, three are needed for the third derivative. \vem

\subproblem{c}
\begin{gather*}
    f'(x)=\frac{3}{x^2}\\
    f(x)=-\frac{3}{x}+C\\
    f(1)=4 \implies C=7\\
    \boxed{f(x)=-\frac{3}{x}+7}
\end{gather*}
The answer is valid on $(0,\infty)$. For $x<0$, as no initial condition is provided on this interval, no further progress can be made.\vem

\subproblem{d}
\begin{gather*}
    f'(x)=1-\frac{1}{x^2}\\
    f(x)=x+\frac{1}{x}+C\\
    f(1)=0 \implies C=-2\\
    \boxed{f(x)=x+\frac{1}{x}-2}
\end{gather*}
The three graphs have the same shape and slope, but they are vertical translations of each other. The theorem that guarantees this is theorem 3 (families of antiderivatives).

\problem{5}

\subproblem{a}
\begin{gather*}
    v(t)=0.6t-0.06t^2\\
    v(5)=1.5m/s\\
    s(t)=0.3t^2-0.02t^3\\
    s(5)=5m
\end{gather*}
The elevator is speeding up on $t\in [0,5]$. \vem

\subproblem{b}
\begin{gather*}
    v(t)=-8t+24\\
    s(t)=-4t^2+24t\\
    v(t)=0 \implies t=3s\\
    s(3)=36m\\
    d=\frac{v_t^2-v_i^2}{2a}=\frac{v_o^2}{2k}
\end{gather*}

\subproblem{c}
\begin{gather*}
    v(t)=-9.8t+14\\
    s(t)=-4.9t^2+14t+2\\
    v(t)=0 \implies t=\frac{10}{7}\\
    s(\frac{10}{7})=12m\\
    s(t)=0 \implies t=2.99
\end{gather*}

\subproblem{d}
\begin{gather*}
    V(t)=\int r(t) \, dt=-t^2+40t+C\\
    t(0)=150 \implies C=150\\
    V(t)=-t^2+40t+150\\
    V(20)-V(0)=400=A
\end{gather*}
The net change equals the area under the curve.

\problem{6}

\subproblem{a}
\begin{gather*}
    \deriv (\sin ^2 x+C)=2 \sin x \cos x\\
    \deriv (-\cos ^2 x+C)=2 \sin x \cos x\\
    \deriv (-\frac{1}{2}\cos 2x+C)=2 \sin x \cos x\\
    -1, -\frac{1}{2}, \frac{1}{2}
\end{gather*}
It is the theorem of equal derivatives. \vem

\subproblem{b}
\begin{gather*}
    \ln(x^2+1)+C=\ln 3+\ln(x^2+1)+C=\ln(3x^2+3)+C\\
    \ln 3\ne 0
\end{gather*}
As $C$ may represent any real value, the first $C$ and second $C$ simply have a difference of $\ln 3$, which absorbs the difference.\vem

\subproblem{c}
\begin{enumerate}[label=(\roman*)]
    \item True, the difference of two constants is a constant.
    \item True, the constant $C$ may vary.
    \item False, it obeys the product rule.
    \item False, missing constant $C$.
\end{enumerate}

\subproblem{d}
The answers agree. The chain rule method for the case here works, as it only creates a constant term 2, while for $(x^2+3)^5$ a linear term is added. 

\problem{7}

\subproblem{a}
\begin{gather*}
    F'(0^+)=\frac{\frac{1}{2}h^2}{h}=0\\
    F'(0^-)=-\frac{\frac{1}{2}h^2}{h}=0\\
    F'(0^+)=F'(0^-)=0=f(0)\\
    F'(x)=|x|=f(x)
\end{gather*}

\subproblem{b}
\begin{gather*}
    x>0 \implies \ddx \ln x=\frac{1}{x}\\
    x<0\implies \ddx \ln-x=-\frac{1}{x}\\
    \ddx \ln |x|=\frac{1}{x}\\
    G(x)=\begin{cases}
        \ln x + 1,\quad x>0\\
        \ln (-x),\quad x<0
    \end{cases}\\
    F(x)=\begin{cases}
        \ln x+C_1,\quad x>0\\
        \ln (-x)+C_2,\quad x<0
    \end{cases}
\end{gather*}
Theorem 2 fails as $F(x)$ is not differentiable on the entire interval.\vem

\subproblem{c}
\begin{gather*}
    x\in (-1,0] \implies f(x)=0\implies F(x)=c\\
    x\in (0,1)\implies f(x)=1\implies F(x)=x+k\\
    \lim_{x \to 0^-} c=\lim_{x \to 0^+} \implies c=k\\
    F'(0^+)=\frac{h}{h}=1\\
    F'(0^-)=\frac{0}{h}=0\\
    F'(0^+)\ne F'(0^-)\\
    F'(x)\ne f(x)
\end{gather*}

\subproblem{d}
A function with a jump discontinuity skips values, and thus failing to meet the requirement of taking every value on an interval as a derivative. A step function is integrable, yet it does not have a derivative. Thus, having an anti-derivative is the stronger condition.















\end{document}

